\section{Experiments}
\label{sec:experiments}

\subsection{Base recipe}
All runs use the MMAudio \texttt{small\_16k} configuration:
latent dim 20, hidden dim 448, 7 attention heads, 12 joint
MMDiT blocks followed by 8 fused blocks, 8\,s clips at 16\,kHz.
We train for 300{,}000 iterations with batch size 512 (summed
over GPUs, per-GPU batch $512 / W$), AdamW ($\beta =
[0.9, 0.95]$, weight decay $10^{-6}$), peak lr $10^{-4}$ with
linear warmup over 1000 steps and step decay by $0.1$ at
240k and 270k, gradient clip $1.0$, bfloat16 autocast (except
inside the GW block), and Post-Hoc EMA \citep{cheng2024mmaudio}
with $\sigma_{\text{rel}} = [0.05, 0.1]$ checkpointed every 5000
steps. Final evaluation uses the synthesized EMA weights at
$\sigma = 0.05$. Every job script in \texttt{jobs/} launches
\texttt{train.py} with exactly these defaults and overrides only
the axis under test.

\subsection{Data}
Training data is the MMAudio multi-corpus mix defined in
\texttt{mmaudio/data/data\_setup.py}, itself the release-time
instantiation of the MMAudio paper's joint-training recipe
(roughly $1{:}1$ audio--visual to audio--text sampling
\citep{cheng2024mmaudio}): VGGSound \citep{chen2024vggsound}
precomputed latents oversampled $5\times$, plus AudioCaps
\citep{kim2019audiocaps}, AudioSetSL, BBCSound, FreeSound, and
Clotho as audio-only streams from precomputed memmap
tensordicts. We do not modify the mix; the relational term in
Sec.\,\ref{sec:method:reg} is the only training change relative
to MMAudio. Features are
precomputed once with
\texttt{training/extract\_video\_training\_latents.py}: VAE
latent $(250, 20)$, CLIP features $(64, 1024)$, Synchformer
features $(192, 768)$, and CLIP-text features $(77, 1024)$.
Validation and evaluation use the VGGSound val and test splits.
The OOD axis uses the 30-class held-out splits produced by
\texttt{experiments/make\_ood\_split.py}.

\subsection{Metrics}
We report four metrics per run via \texttt{av-benchmark}
(\texttt{mmaudio/runner.py::eval}):
\begin{description}[leftmargin=2em]
  \item[FD$_{\text{PaSST}}$] Fr\'echet distance in a PaSST audio
    embedding space. Lower is better; measures distribution-level
    fidelity.
  \item[IS] Inception Score on PaSST logits. Higher is better;
    measures diversity and per-sample recognizability.
  \item[IB-score] ImageBind
    \citep{girdhar2023imagebind} cross-modal retrieval score
    between generated audio and the video. Higher is better;
    measures semantic coupling.
  \item[DeSync] Synchformer-based temporal misalignment score.
    Lower is better; measures timing.
\end{description}
No single metric dominates --- in particular, IS and
FD$_{\text{PaSST}}$ measure overlapping but not identical things,
and IB-score and DeSync can trade off. The full row of four is
what we compare.

\subsection{Hypotheses}
\label{sec:experiments:hypotheses}
\begin{description}[leftmargin=2em]
  \item[H1 (does it help at all?)]
    The best GW variant at default $\lambda_{\mathrm{GW}}{=}0.01$
    improves over the no-GW baseline on at least 2 of 4 metrics
    without regressing on the remaining two by more than 5\%
    relative. \emph{Tested by:} baseline vs D1 variant sweep.
  \item[H2 (representation matters)]
    The four D1 variants produce materially different outcomes
    (std of any one metric across variants exceeds 2 $\times$ the
    baseline-to-best gap), so the choice of representation is
    a real design decision and not noise.
    \emph{Tested by:} D1 variant sweep.
  \item[H3 ($\lambda$ has an optimum)]
    The metric-vs-$\lambda$ curve is non-monotone at the best
    variant: too-small $\lambda$ is indistinguishable from
    baseline, too-large $\lambda$ degrades FD$_{\text{PaSST}}$ by
    pulling the audio geometry away from the data manifold.
    \emph{Tested by:} D2 lambda sweep (\texttt{BEST\_VARIANT}).
  \item[H4 (direction is asymmetric)]
    \texttt{detach\_video=True} outperforms
    \texttt{detach\_video=False} on IB-score and
    FD$_{\text{PaSST}}$; letting the video side adapt creates
    shortcuts that improve $\LGW$ at the cost of
    representational quality.
    \emph{Tested by:} D3 detach ablation.
  \item[H5 (warmup matters; cosine decay late)]
    A late-game schedule (\texttt{cosine\_anneal}) matches or
    beats \texttt{constant} on fidelity metrics
    (FD$_{\text{PaSST}}$) because the final iterations are free
    to polish under pure flow matching; \texttt{linear\_rampup}
    is no worse than \texttt{constant}.
    \emph{Tested by:} D4 schedule ablation.
\end{description}

\subsection{Ablation grid}
\label{sec:experiments:grid}

\begin{table}[t]
\centering
\small
\begin{tabular}{@{}llll@{}}
\toprule
Job script & Axis & Runs & Core override \\
\midrule
\texttt{train\_baseline.job} & (none)   & 1 & \texttt{gw\_regularization.enabled=false} \\
\texttt{train\_variants.job} & D1       & 4 & \texttt{gw\_regularization.variant=\{global,projected,c\_g,fused\}} \\
\texttt{train\_lambda.job}   & D2       & 5 & \texttt{gw\_regularization.lambda\_gw=\{0.001,...,0.1\}} \\
\texttt{train\_detach.job}   & D3       & 2 & \texttt{gw\_regularization.detach\_video=\{true,false\}} \\
\texttt{train\_schedule.job} & D4       & 3 & \texttt{gw\_regularization.schedule=\{constant,linear\_rampup,cosine\_anneal\}} \\
\texttt{train\_ood.job}      & OOD      & 2 & \texttt{data.\{VGGSound,ExtractedVGG\}=<ood tsv>} \\
\texttt{train\_gwxsync.job}  & D5       & 4 & $2\times 2$ GW $\times$ \texttt{disable\_sync} \\
\bottomrule
\end{tabular}
\caption{Training runs. 1 baseline + 20 treatment runs. All
eval runs go through \texttt{jobs/eval.job} (or
\texttt{eval\_all.job} for the array version), which loads the
EMA-final weights and calls \texttt{batch\_eval.py} followed by
\texttt{av-benchmark}.}
\label{tab:grid}
\end{table}

We run the grid in two phases. Phase 1: baseline + D1 variant
sweep. Phase 2 (conditional on the winning variant
\texttt{BEST\_VARIANT}): D2 / D3 / D4 / D5 plus OOD.

\subsection{Analysis beyond aggregate metrics}
\label{sec:experiments:analysis}
We complement FD/IS/IB/DeSync with two diagnostics written as
\texttt{experiments/analysis.py} subcommands:
\begin{description}[leftmargin=2em]
  \item[\texttt{geometry}] Effective rank
    (exp--entropy of normalized singular spectrum) and the full
    normalized spectrum of the projected video and audio
    representations on a held-out batch. GW is predicted to
    \emph{raise} effective rank by preventing collapse to a
    pointwise subspace.
  \item[\texttt{perclass}] Per-class IB-score differences
    (GW $-$ baseline) on VGGSound test, shown as a histogram and
    a ranked CSV. If GW aligns relational structure rather than
    just shifting an aggregate, the distribution of deltas should
    be broad, not a single spike at the mean.
\end{description}
\texttt{experiments/analyze\_couplings.py} additionally visualizes
the optimal couplings $\Tstar$ grouped by class pair: a heatmap
over the 20 most frequent VGGSound classes, plus a ranked CSV of
per-class self-alignment (the average within-class diagonal
mass). A strong GW variant should show heavier within-class mass
and clean block-diagonal structure; a collapsed one should show
uniform mass.

\subsection{Runtime and reproducibility}
Each 300k-step run fits on a single GPU under 48 hours on
\texttt{acltr}. Slurm headers are fixed across jobs via
\texttt{jobs/\_header.sh} and \texttt{jobs/\_common.sh}: CUDA
12.1.1 + NCCL 2.18.3, conda env \texttt{gwot}, thread pools
explicitly capped (\texttt{OMP\_NUM\_THREADS=4} etc.) to stay
below the cluster's per-user \texttt{RLIMIT\_NPROC}. Re-running
the same \texttt{exp\_id} auto-resumes from
\texttt{output/<exp\_id>/<exp\_id>\_ckpt\_last.pth}, so a
pre-empted run picks up where it left off.
